%==============================================================================
\section{Phân cụm dựa trên mật độ (Density-Based Clustering)}
\label{sec:density-based}
%==============================================================================

\begin{dinhnghia}[title={Ý tưởng phân cụm theo mật độ}]
Phân cụm dựa trên mật độ cho rằng cụm là các vùng \textbf{mật độ cao}, ngăn cách bởi vùng \textbf{mật độ thấp}.
\end{dinhnghia}

\begin{phuongphap}[title={Ba loại điểm}]
\begin{enumerate}
  \item \textbf{Core (lõi)}: điểm có ít nhất một số láng giềng tối thiểu trong bán kính cho trước --- tâm vùng cụm.
  \item \textbf{Border (biên)}: không phải core nhưng nằm trong lân cận của ít nhất một điểm core.
  \item \textbf{Noise (nhiễu)}: không phải core cũng không phải border.
\end{enumerate}
\end{phuongphap}

\begin{ynghia}[title={So với centroid/partition}]
Không giả định cụm hình cầu hay số cụm cố định; phù hợp hình dạng bất kỳ và phát hiện điểm nhiễu.
\end{ynghia}

\subsection{Các thuật toán phân cụm theo mật độ phổ biến}

\subsubsection{DBSCAN}

\begin{dinhnghia}[title={DBSCAN}]
DBSCAN (Density-Based Spatial Clustering of Applications with Noise) là thuật toán phân cụm theo mật độ phổ biến nhất. Cần hai tham số:
\begin{itemize}
  \item \texttt{minPts}: số láng giềng tối thiểu để một điểm được coi là core.
  \item \texttt{eps}: bán kính (hoặc khoảng cách tối đa) xét láng giềng giữa các điểm core.
\end{itemize}
\end{dinhnghia}

\begin{phuongphap}[title={Cơ chế tóm tắt}]
Từ mỗi core chưa thăm, mở rộng cụm bằng cách nối các core và border liên thông trong $\varepsilon$; điểm không đạt điều kiện gán là noise.
\end{phuongphap}

\subsubsection{OPTICS}

\begin{dinhnghia}[title={OPTICS}]
OPTICS (Ordering Points To Identify the Clustering Structure) xây \textbf{đồ thị reachability} nối mỗi điểm tới láng giềng gần trong ngưỡng khoảng cách; cạnh có trọng số theo khoảng cách. Từ đồ thị này, thuật toán tạo cấu trúc phân cấp bằng cách tách theo ngưỡng mật độ.
\end{dinhnghia}

\begin{tinhchat}[title={So với DBSCAN}]
OPTICS xử lý tốt hơn khi cụm có \textbf{mật độ khác nhau}, có nhiễu, và khi cần xem cấu trúc phân cấp mật độ mà không cố định một cặp \texttt{eps}/\texttt{minPts} duy nhất.
\end{tinhchat}

\subsubsection{HDBSCAN}

\begin{dinhnghia}[title={HDBSCAN}]
HDBSCAN (Hierarchical Density-Based Spatial Clustering of Applications with Noise) mở rộng DBSCAN theo hướng phân cấp mật độ.
\end{dinhnghia}

\begin{tinhchat}[title={Ưu điểm so với DBSCAN}]
\begin{itemize}
  \item Cụm \textbf{mật độ thay đổi} trong cùng dataset.
  \item Phát hiện cụm \textbf{hình dạng và kích thước đa dạng}.
  \item Ít phụ thuộc chọn một $\varepsilon$ duy nhất cho mọi mức mật độ.
\end{itemize}
\end{tinhchat}

\begin{luuy}[title={Các chương tiếp theo}]
Ba chương kế tiếp trình bày chi tiết DBSCAN, OPTICS và HDBSCAN kèm triển khai Python.
\end{luuy}
